\documentclass[a4paper, 11pt]{article}

\usepackage{revy}
\usepackage[T1]{fontenc}
\usepackage{amsfonts}
\usepackage{amssymb}
\usepackage[utf8]{inputenc}
% Required for inserting images

\revyname{Matematikrevyen}
\revyyear{2025}
\version{0.1}
\eta{$3$ minutter}
\status{ Ikke færdig}

\title{Matematikkens historie}
\author{Malthe '20, Christian O '24, Louie '24, Sirius '23}

\begin{document}

\maketitle

\begin{roles}  
\role{I}[KE] Instruktør
\role{F}[Niklas] En forelæser
\role{S1}[Conrad] En studerende
\role{S2}[Alberte F] En studerende
\role{St1}[Freja] En studerende (statist)
\role{St2}[Markus] En studerende (statist)
\role{St3}[Oliver] En studerende (statist)
\end{roles}

\begin{sketch}
\scene Man prøver at bilde nogen ind at alle matematiske begreber er opkaldt efter folk. Forelæser står med en computer og trykker sig igennem slides, som vises på AV. Slidesne viser sort-hvid billeder af de nævnte personer og deres navne.

\says{F} Hej, og velkommen til matematikkens historie. I dag skal vi lære om oprindelsen af forskellige matematiske begreber. Vores første eksempel er "Algebra". Selve emnet har rødder tilbage til antikken, men navnet kommer faktisk fra det 9. århundrede og er opkaldt efter Al-Khwarizmis bog "Al-Jabr". 

\says{S1}[til S2] Wow, det er faktisk en sjov fact.   

\says{F} Ja, ikke! Så tror jeg også I vil synes godt om den næste. Analyse er jo faktisk opkaldt efter "Anna de Lucia" fra Andalusien, der var en af de første matematikere, der formaliserede funktionsbegrebet.

\says{S2}[til S1] Okay, det havde jeg faktisk ikke regnet med.

\says{F} "Geometri" er jo faktisk opkaldt efter den russiske matematiker "Grigorij Gimetrij", der tegnede de første paralelle linjer tilbage i 1743. 

\says{S1}[til S2] Jeg troede at "geometri" kom fra græsk?

\says{F} Shh, du må række hånden op, hvis du har noget at sige. Vores næste eksempel er den hollandske matematiker "Woude Statistijk", der grundlagde statistikken for at vise, at han ikke var gennemsnitlig. 

\says{S2}[til S1] Ej, nu finder han bare på ting.

\says{F} Og vi må jo ikke glemme russeren "Rasputin Taltjerri", grundlæggeren af moderne talteori. 

\says{S1}[Højlydt] Du kan da ikke mene, at talteori blev opfundet af en russisk matematiker ved navn "Taltjerri"!

\says{F} Nej, det var jo heller ikke det jeg sagde. Han grundlagde den \textit{moderne} talteori. Selve ideen om tal går jo helt tilbage til verdensmesteren i skak i 1960, den lettisk-sovjetiske "Mikhail Tal".

\says{S2} Stop! Hvis du påstår at én til åndssvag ting er opkaldt efter en matematiker, så går jeg min vej. 

\says{F} Hov hov du, nu slapper du lige lidt af. 

\scene AV skifter til et billede af Johan Jensen med en baret og et overskæg. Titlen er "Jensens ulighed".

\says{F} Vi har selvfølgelig også den vigtige sætning om konkave funktioner, Jensens ulighed. Den er selvfølgelig opkaldt efter den franske matematiker "François Jensôn". 

\says{S2} That's it, jeg går min vej.

\says{F} Jeg har kun et enkelt slide tilbage, bliv lige hængende.

\scene AV viser et billede af Neo og Morpheus fra "The Matrix".

\says{S2} Lad mig gætte. Du vil prøve at bilde mig ind at ordet matrix blev opfundet i 1999? 

\says{F} Nej, selvfølgelig ikke, lineær algebra har rødder tilbage til Descartes, Leibniz og Grassmann. Det er gruppeteori der fik et væsentligt bidrag, da filmen introducerede konceptet "isomorpheus".  

\scene Lys ned
\end{sketch}

\end{document}